\section{Introduction}
\label{sec:introduction}

Federated learning (FL) trains shared models without pooling raw traffic~\cite{mcmahan2017communication}.
Reconstruction-error autoencoders (AEs) are a natural fit for IoT intrusion detection
because benign-only training suffices and attack labels are not required at training time.

A standard deployment trains one shared AE via FedAvg, then applies a threshold $\tau$
above which a reconstruction error signals an intrusion.
In homogeneous fleets this works well.
In heterogeneous fleets, however, benign reconstruction-error distributions differ
substantially across device types: a single shared $\tau$ imposes unequal false-positive
rates (FPRs).
Some devices carry a disproportionate false-alarm burden---making them noisy
or operationally unusable---while others may become under-sensitive.
A shared server-broadcast threshold is natural in deployments where the system
exposes a single global alarm boundary or where per-client threshold state is
not managed independently.

This study asks how large the per-client FPR disparity is under a shared threshold,
how much per-client calibration reduces it, and what detection-quality tradeoff
this creates---all under a fixed FedAvg AE with real physical-device heterogeneity.
The target operational scenario is a fleet of physical IoT devices for which
local benign calibration data are available, and where reducing unequal
false-alarm burden across devices is a deployment priority.
Prior FL-AE work may report aggregate detection metrics or use thresholds implicitly,
but without fixing the encoder, aggregation protocol, and training hyperparameters and
varying only threshold scope, it is impossible to attribute per-client alarm-rate
dispersion to the threshold policy rather than to training differences.
This controlled-comparison gap motivates the present study.

The primary metric is the coefficient of variation of FPR across eligible clients,
$\text{CV(FPR)} = \sigma_{\text{FPR}} / \mu_{\text{FPR}}$.
Four threshold policies are evaluated: a client-averaged shared threshold (B1), full per-client
personalization (B2), device-family grouping (B3), and unsupervised cluster grouping (B4).
B1 is constructed as the arithmetic mean of eligible clients' local 95th-percentile
benign calibration thresholds, broadcast as a single scalar.
The arithmetic mean treats eligible clients equally and models a simple server-broadcast
threshold derived from local calibration summaries, avoiding sample-count domination;
pooled and sample-weighted alternatives are reported as sensitivity checks.
A centralized baseline (B0) is included as a privacy-incompatible reference
comparator but is not part of the single-variable B1--B4 comparison.

The study uses N-BaIoT~\cite{meidan2018n} as the primary confirmatory regime (Regime~A),
CICIoT2023~\cite{neto2023ciciot2023} as a file-level applicability check under
near-homogeneous conditions (Regime~B), and synthetic Dirichlet partitions of N-BaIoT as a
heterogeneity severity sweep (Regime~C).
Five random seeds are used; seed-level descriptive summaries (mean, SD, range, sign consistency) report stability.
The FL simulation uses Flower~\cite{beutel2020flower}.

\textbf{Contributions.}
\begin{itemize}
  \item \emph{Scientific contribution:} A controlled threshold-calibration study
        fixing the FedAvg AE ($E{=}1$, full participation), training protocol, and seeds
        across B1--B4; only the post-training threshold scope varies.
        Because B2 uses each client's own benign $p_{95}$ threshold, FPR equalization
        is expected by construction under stable distributions; the contribution is the
        controlled held-out measurement, seed consistency, tradeoff characterization,
        and B4 approximation---all under a fixed encoder.
  \item \emph{Measured result:} B2 reduces CV(FPR) from $\approx 1.02$ to $0.299$
        (95\% bootstrap CI on $\Delta$: $[0.647,\,0.769]$, all five seed deltas positive;
        relative reduction $(\text{CV}_{\text{B1}} - \text{CV}_{\text{B2}})/\text{CV}_{\text{B1}} = 70.6\%$).
        The detection-quality cost is characterized: P10 Macro-F1 degrades in a single
        low-separability client.
  \item \emph{Scope and limits:} Under the CICIoT2023 randomized file-level pseudo-client
        partition (not physical devices), no improvement is observed.
        A Dirichlet severity sweep is consistent with the benefit concentrating under
        stronger heterogeneity.
  \item \emph{Replication artifact:} Source code, configs, per-seed results, and
        table/figure generation scripts archived on Zenodo
        (DOI: \nolinkurl{10.5281/zenodo.20315654}).
\end{itemize}
